% 引力波天文学（综述）
% license CCBYSA3
% type Wiki

本文根据 CC-BY-SA 协议转载翻译自维基百科\href{https://en.wikipedia.org/wiki/Gravitational-wave_astronomy}{相关文章}。


引力波天文学是天文学的一个分支，主要研究由天体物理源发出的引力波的探测与分析。\(^\text{[1]}\)

引力波是时空中极其微小的扭曲或涟漪，由大质量天体的加速运动所引起。它们产生于灾变性事件之中，例如双黑洞并合、双中子星合并、超新星爆炸，以及大爆炸后早期宇宙中的物理过程。研究引力波为观测宇宙提供了一种全新的方式，使人们能够深入了解物质在极端条件下的行为。类似于电磁辐射（如光波、无线电波、红外辐射和 X 射线）通过电磁场涨落传播能量的方式，引力辐射则是由相对较弱的引力场涨落所产生的能量传播。引力波的存在最早由奥利弗·赫维赛德（Oliver~Heaviside）于 1893 年提出，随后在 1905 年由昂利·庞加莱（Henri~Poincaré）作为电磁波的引力学对应物加以推测，最终在 1916 年被阿尔伯特·爱因斯坦（Albert~Einstein）根据广义相对论所预言。

1978 年，罗素·阿伦·赫尔斯（Russell~Alan~Hulse）与约瑟夫·胡顿·泰勒（Joseph~Hooton~Taylor~Jr.）通过观测两颗互相绕转的中子星，首次提供了引力波存在的实验性证据。二人因此工作获得了 1993 年诺贝尔物理学奖。2015 年，距爱因斯坦的预言近一个世纪后，人类首次直接探测到了由双黑洞并合产生的引力波信号，从而确认了这一长期难以捉摸的现象的存在，并开启了天文学的新纪元。此后，科学家又陆续探测到多起包括双黑洞并合、中子星碰撞及其他剧烈宇宙事件在内的引力波信号。如今，引力波的探测主要依靠激光干涉测量技术（laser interferometry），该技术通过测量两条相互垂直臂长随波动变化的微小差异来捕捉引力波。诸如 LIGO（Laser~Interferometer~Gravitational-wave~Observatory，激光干涉引力波天文台）、Virgo 与 KAGRA（Kamioka~Gravitational~Wave~Detector，神冈引力波探测器）等观测站都采用此项技术来探测来自遥远宇宙事件的微弱信号。LIGO 的联合创始人巴里·C·巴里什（Barry~C.~Barish）、基普·S·索恩（Kip~S.~Thorne）与赖纳·魏斯（Rainer~Weiss）因其在引力波天文学领域的开创性贡献而共同获得了 2017 年诺贝尔物理学奖。

当利用电磁波观测遥远的天体时，散射、吸收、反射、折射等多种物理过程会导致信息丢失。  
宇宙中存在若干光子只能部分穿透的区域，例如星云内部、银河系核心的致密尘埃云以及黑洞附近的区域等。与电磁天文学相互补充的引力波天文学，有潜力以更高的分辨率研究宇宙。在一种被称为“多信使天文学”（multi-messenger astronomy）的方法中，引力波数据与不同波段的电磁观测数据相结合，以获得对天体物理现象更全面的理解。引力波天文学不仅有助于揭示早期宇宙的状态、检验引力理论、探测暗物质与暗能量的分布，还能帮助确定哈勃常数（Hubble constant），即描述宇宙加速膨胀速率的重要参数。这些研究为超越标准模型（Beyond the Standard Model, BSM）的物理学打开了新的大门。


当前该领域仍面临多项挑战，包括噪声干扰、超高灵敏度探测设备的缺乏以及低频引力波信号的探测难题。地面探测器受到环境扰动引起的地震振动影响，同时其探测臂长度也受到地球曲率的限制。未来，引力波天文学将致力于研发升级型探测装置与下一代观测台，同时探索基于太空的探测方案，例如 LISA（Laser~Interferometer~Space~Antenna，激光干涉空间天线）。  
LISA 将能够探测来自银河系核心致密超大质量黑洞、原初黑洞等遥远源的信号，以及诸如双白矮星并合与早期宇宙来源的低频引力波信号。\(^\text{[2]}\)
\subsection{引言}
引力波是由轨道双体系统中加速运动的质量所产生的引力强度扰动波，它们以光速从源处向外传播。这种波动最早由奥利弗·赫维赛德（Oliver~Heaviside）于 1893 年提出，后又由昂利·庞加莱（Henri~Poincaré）于 1905 年进一步提出，认为它们类似于电磁波，但属于引力的对应形式。


1916 年，阿尔伯特·爱因斯坦（Albert~Einstein）在其广义相对论的框架下正式预言了引力波的存在，并将其描述为时空中的涟漪。然而，爱因斯坦后来一度拒绝接受引力波的存在。\(^\text{[3]}\)引力波以引力辐射（gravitational~radiation）的形式传输能量，这种辐射是一种类似于电磁辐射的辐射能。经典力学中的牛顿万有引力定律并不包含引力波的存在，因为该定律假设物理作用是瞬时传播（即传播速度无限大）的——这恰好揭示了牛顿力学在解释相对论相关现象方面的局限性。

引力波存在的第一个间接证据于 1974 年发现，来自对赫尔斯–泰勒双星脉冲星（Hulse–Taylor~binary~pulsar）轨道衰变的观测，其衰减速率与广义相对论预测的因引力辐射能量损失所导致的结果一致。1993 年，罗素·A·赫尔斯（Russell~A.~Hulse）与约瑟夫·胡顿·泰勒（Joseph~Hooton~Taylor~Jr.）因这一发现共同获得诺贝尔物理学奖。

直到 2015 年，人类才首次直接观测到引力波信号——该信号来自两颗黑洞并合事件，被美国路易斯安那州利文斯顿（Livingston,~Louisiana）与华盛顿州汉福德（Hanford,~Washington）的 LIGO（Laser~Interferometer~Gravitational-wave~Observatory，激光干涉引力波天文台）探测器接收到。  
2017 年诺贝尔物理学奖随后授予赖纳·魏斯（Rainer~Weiss）、基普·索恩（Kip~Thorne）与巴里·巴里什（Barry~Barish），以表彰他们在引力波直接探测中的奠基性贡献。

在引力波天文学（gravitational-wave~astronomy）中，科学家通过观测引力波来推断其源的物理性质。  
能够通过这种方式研究的源包括：由白矮星、中子星或黑洞组成的双星系统；超新星爆发事件；以及宇宙大爆炸后早期宇宙形成阶段的物理过程。
\subsubsection{仪器与挑战}
由于不同探测器在技术参数与灵敏度上的差异，探测器间的协作有助于收集互补且极具价值的数据。目前存在多台地面激光干涉仪，长度跨越数英里或数公里，包括：位于美国华盛顿州与路易斯安那州的两座 LIGO 探测器；位于意大利欧洲引力观测台（European~Gravitational~Observatory）的 Virgo 探测器；位于德国的 GEO600；  
以及位于日本的神冈引力波探测器 KAGRA（Kamioka~Gravitational~Wave~Detector）。迄今为止，LIGO、Virgo 与 KAGRA 已开展多次联合观测；而 GEO600 由于其设备灵敏度相对较低，目前主要用于实验性测试与技术验证，近年未参与其他观测台的联合运行。
\begin{figure}[ht]
\centering
\includegraphics[width=10cm]{./figures/e66390bc2706521f.png}
\caption{不同类型引力波探测器的噪声曲线（noise~curves）随频率变化的示意图。在极低频区域为脉冲星计时阵列（Pulsar~Timing~Arrays），在低频区域为空间搭载探测器（space-borne~detectors），而在高频区域为地面探测器（ground-based~detectors）。图中同时显示了潜在天体物理源的特征应变（characteristic~strain）。要使信号可被探测，其特征应变必须高于噪声曲线。\(^\text{[4]}\)} \label{fig_Gras_1}
\end{figure}
\subsubsection{高频引力波}
2015 年，LIGO（Laser~Interferometer~Gravitational-wave~Observatory，激光干涉引力波天文台）项目首次利用激光干涉仪直接探测到引力波。\(^\text{[5][6]}\) 
LIGO 探测器观测到来自两颗恒星质量黑洞并合的引力波信号，其结果与广义相对论的预测完全一致。\(^\text{[7][8][9]}\)这些观测不仅证实了双恒星质量黑洞系统的存在，也实现了人类历史上首次引力波的直接探测与首次双黑洞并合的观测。\(^\text{[10]}\)这一发现被认为是科学上的一次革命性突破，因为它验证了人类利用引力波天文学探索宇宙、研究暗物质与宇宙大爆炸起源的能力。
\subsubsection{低频引力波}
另一种观测手段是利用脉冲星计时阵列（Pulsar~Timing~Arrays,~PTAs）。目前存在三个主要合作组织：欧洲脉冲星计时阵列（European~Pulsar~Timing~Array,~EPTA）、  
北美纳赫兹引力波天文台（North~American~Nanohertz~Observatory~for~Gravitational~Waves,~NANOGrav）以及帕克斯脉冲星计时阵列（Parkes~Pulsar~Timing~Array,~PPTA），  
三者共同组成国际脉冲星计时阵列（International~Pulsar~Timing~Array）。这些项目使用现有的射电望远镜进行观测。由于其灵敏频率处于纳赫兹（nanohertz）范围，因此需要长达数年的持续观测才能探测到引力波信号，且探测灵敏度会随时间逐步提高。当前的灵敏度上限已接近理论上对天体物理源的预期值。\(^\text{[11]}\)
\begin{figure}[ht]
\centering
\includegraphics[width=10cm]{./figures/d7c4d5b188562cf1.png}
\caption{NANOGrav（2023）观测到的脉冲星之间相关性随其角距离变化的关系图，其中虚线紫色曲线表示理论 Hellings--Downs模型的预期结果，实线绿色曲线表示在不存在引力波背景情况下的相关性。\(^\text{[12][13]}\)} \label{fig_Gras_2}
\end{figure}
2023 年 6 月，四个脉冲星计时阵列（Pulsar~Timing~Array,~PTA）合作组织——前述的三个（EPTA、NANOGrav 与 PPTA）以及中国脉冲星计时阵列（Chinese~Pulsar~Timing~Array,~CPTA）—— 分别独立发布了相似的观测结果，均表明存在一个纳赫兹（nanohertz）频率范围内的随机引力波背景（stochastic~background~of~gravitational~waves）。\(^\text{[14]}\)每个团队都独立地首次测量到了理论预期的Hellings--Downs 曲线，即两个脉冲星在天空中角距离的函数形式下的四极相关性（quadrupolar~correlation），这是观测到的背景辐射具有引力波起源的典型标志。\(^\text{[15][16][17][18]}\)该背景辐射的具体来源尚未确定，但超大质量黑洞双星系统（binaries~of~supermassive~black~holes）被认为是最有可能的候选者。\(^\text{[19]}\)
\subsubsection{中频引力波}
未来，引力波观测还可能通过太空搭载的探测器得以扩展。欧洲航天局（European~Space~Agency,~ESA）已将一项引力波任务选定为其 L3 计划的核心任务，计划于 2034 年发射。当前的任务方案为进化版激光干涉空间天线（evolved~Laser~Interferometer~Space~Antenna,~eLISA）。\(^\text{[20]}\)与此同时，日本也在研制分赫兹干涉式引力波天文台（Deci-hertz~Interferometer~Gravitational-wave~Observatory,~DECIGO），  
旨在探测介于地面与太空观测器之间频率范围的引力波信号。
\subsection{科学价值}
天文学传统上依赖于电磁辐射进行观测。最初仅限于可见光波段，随着技术的进步，逐渐能够观测到从射电波（radio~waves）到伽马射线（gamma~rays）的电磁波谱各个部分。每当新的频段被探测，人类便获得了观察宇宙的新视角，并由此迎来新的发现。\(^\text{[21]}\)20 世纪期间，对高能量和高质量粒子的间接测量以及后来的直接测量，为探索宇宙提供了另一扇窗口。20 世纪晚期，太阳中微子的探测奠定了中微子天文学（neutrino~astronomy）的基础，使人类得以窥见此前无法接近的现象，例如太阳的内部机制。\(^\text{[22][23]}\)引力波的观测为进行天体物理研究提供了又一种新的途径。


罗素·赫尔斯（Russell~Hulse）与约瑟夫·泰勒（Joseph~Taylor）因观测到一对中子星（其中之一为脉冲星）的轨道衰变符合广义相对论关于引力辐射的预言,而获得了 1993 年诺贝尔物理学奖。\(^\text{[24]}\)此后，又发现了许多其他双星脉冲星系统（包括一个双脉冲星系统），
其观测结果均与引力波理论预测一致。\(^\text{[25]}\)2017 年，瑞纳·魏斯（Rainer~Weiss）、基普·索恩（Kip~Thorne）与巴里·巴里什（Barry~Barish）因在首次直接探测引力波中的关键贡献而获得诺贝尔物理学奖。\(^\text{[26][27][28]}\)

引力波为天文学提供了与其他观测手段互补的信息。通过将同一事件的多种观测结果结合（如电磁波、引力波、中微子信号等），可以更全面地理解天体源的物理性质——这种方法被称为 多信使天文学（multi-messenger~astronomy）。此外，引力波还可用于研究那些在其他波段几乎不可见甚至完全无法探测的系统。例如，引力波为测量黑洞性质提供了一种独特而直接的途径。


引力波可由多种系统发射，但要产生可探测的信号，其源必须由极其巨大的天体组成，并以接近光速的显著比例运动。主要的引力波源是由两个致密天体（compact~objects）构成的双星系统。
典型系统包括：
\begin{itemize}
  \item 致密双星系统（Compact~binaries）：由两颗紧密轨道运行的恒星质量天体组成，例如白矮星（white~dwarfs）、中子星（neutron~stars）或黑洞（black~holes）。轨道较宽、频率较低的双星系统是空间探测器（如 LISA）的信号源。\(^\text{[29][30]}\)轨道较近的双星系统则会在地面探测器（如 LIGO）中产生信号。\(^\text{[31]}\)地面探测器还可能探测到包含中等质量黑洞（约数百个太阳质量）的双星系统。\(^\text{[32][33]}\)
  \item 超大质量黑洞双星（Supermassive~black~hole~binaries）：由质量在 $10^{5}$--$10^{9}$ 个太阳质量之间的两颗黑洞组成。超大质量黑洞通常位于星系中心，当星系并合时，其中心黑洞也预期会随之并合。\(^\text{[34]}\)这些系统可能是宇宙中最强的引力波信号源。最大质量的此类双星是脉冲星计时阵列（PTA）的探测对象。\(^\text{[35]}\)较低质量（约一百万太阳质量）的系统则是空间探测器（如 LISA）的信号源。\(^\text{[36]}\)
  \item 极端质量比系统（Extreme-mass-ratio~systems）：指一颗恒星质量的致密天体绕超大质量黑洞运行的系统。\(^\text{[37]}\)它们是 LISA 等空间探测器的重要观测目标。\(^\text{[36]}\)轨道高度偏心的系统在近日点附近会产生一次引力波爆发（burst）；\(^\text{[38]}\)轨道接近圆形的系统（通常出现在并合晚期）会在 LISA 频率带内连续辐射。\(^\text{[39]}\)极端质量比螺旋并合（EMRIs）可在多个轨道周期内被观测，因而成为研究背景时空几何、精确检验广义相对论的理想探针。c
\end{itemize}
除双星系统外，还存在其他潜在的引力波源：
\begin{itemize}
  \item 超新星（Supernovae）：爆发过程中会产生高频引力波爆发，
  可被 LIGO 或 Virgo 等地面探测器探测。\(^\text{[41]}\)
  \item 旋转中子星（Rotating~neutron~stars）}：若存在轴对称破缺，
  将连续辐射高频引力波。\(^\text{[42][43]}\)
  \item 早期宇宙过程（Early~universe~processes）：
  包括暴胀（inflation）或相变（phase~transition）等。\(^\text{[44]}\)
  \item 宇宙弦（Cosmic~strings）：若其存在，也可能发射引力辐射。\(^\text{[45]}\)
  对这些引力波的探测将验证宇宙弦的存在。
\end{itemize}
引力波与物质的相互作用极弱，这既使其难以探测，也意味着它们能在宇宙中自由传播，不会像电磁辐射那样被吸收或散射。因此，我们得以直接“看到”致密系统的核心区域，例如超新星爆心或银河系中心。此外，引力波还可让我们观察到比电磁辐射更早的宇宙时期—— 因为在复合（recombination）之前，早期宇宙对光是不透明的，但对引力波却是透明的。\(^\text{[46]}\)

引力波可自由穿过物质，这也意味着引力波探测器不同于传统望远镜，并非指向单一视场，而是同时“观察”整个天空。探测器在某些方向上更灵敏，因此建立由多个探测器组成的网络具有显著优势。\(^\text{[47]}\)然而，由于当前探测器数量有限，其方向定位能力仍较为有限。
\subsubsection{宇宙暴胀中的引力波}
宇宙暴胀（cosmic~inflation）是一个假设时期，指宇宙在大爆炸（Big~Bang）后的最初约 $10^{-36}$ 秒内经历的极其迅速的膨胀阶段。这一暴胀过程被认为会产生引力波（gravitational~waves），并在宇宙微波背景辐射（CMB,~cosmic~microwave~background）的偏振模式中留下特征性印记。\(^\text{[48][49]}\)

通过测量微波辐射图样，可以推算原初引力波（primordial~gravitational~waves）的物理性质，并利用这些计算结果来研究早期宇宙的状态与演化。
\subsection{发展}
\begin{figure}[ht]
\centering
\includegraphics[width=6cm]{./figures/750028016b808344.png}
\caption{LIGO 汉福德控制室} \label{fig_Gras_4}
\end{figure}
作为一门仍处于起步阶段的研究领域，引力波天文学（gravitational\mbox{-}wave~astronomy）仍在持续发展中；然而，在天体物理学界已形成共识：这一领域将演化为二十一世纪多信使天文学（multi\mbox{-}messenger~astronomy）的核心组成部分之一。\(^\text{[50]}\)

引力波观测可补充电磁波谱（electromagnetic~spectrum）范围内的观测结果。\(^\text{[50][51]}\)同时，引力波还提供了一种全新的途径，能以传统电磁波探测和分析所无法实现的方式获取信息。电磁波可能被吸收或再辐射，使得对源信息的提取变得困难；而引力波与物质相互作用极弱，不会被散射或吸收。这使得天文学家能够以前所未有的方式观测超新星核心、恒星星云、乃至正在碰撞的星系核心。

地基探测器（ground\mbox{-}based~detectors）已经为恒星级黑洞双星系统的螺旋并合阶段（inspiral~phase~and~merger）、以及中子星双星的并合提供了新的信息。它们还可能探测来自核心坍缩超新星（core\mbox{-}collapse~supernovae）或周期性源（如轻微形变的脉冲星）的信号。若关于早期宇宙（宇宙时间约为 $10^{-25}$ 秒）中某些相变（phase~transitions）或长宇宙弦（cosmic~strings）产生的“弯结爆发”（kink~bursts）的推测属实，这些信号也有可能被探测到。\(^\text{[52]}\)

太空探测器（space\mbox{-}based~detectors），如激光干涉空间天线
（LISA,~Laser~Interferometer~Space~Antenna），应当能够探测由两颗白矮星组成的双星系统、以及 AM~CVn 型星（即一颗白矮星从其低质量氦伴星吸积物质）。此外，LISA 还可观测超大质量黑洞（supermassive~black~holes）的并合，以及小质量天体（质量约为 $1$ 至 $10^3$ 个太阳质量）向此类黑洞的螺旋吸入（inspiral）。LISA 还将能够探测与地基探测器相同类型的早期宇宙信号，但在更低频段下、以更高灵敏度进行观测。\(^\text{[53]}\)

探测引力波是一项极其艰难的任务。这需要超高稳定度的高品质激光器，以及经精密校准的探测器，其灵敏度至少应达到$2\!\times\!10^{-22}\ \mathrm{Hz}^{-1/2}$，如地基探测器 GEO600 所示。\(^\text{[54]}\)研究还指出，即使来自超新星爆炸等巨大的天文事件，这些波在传播过程中也可能衰减为仅相当于原子直径的微小振动。\(^\text{[55]}\)

确定引力波来源的位置同样充满挑战。然而，通过引力透镜（gravitational~lensing）造成的波前偏折，结合机器学习（machine~learning）分析，有望使定位更加容易且精确\(^\text{[56]}\)正如 SN~Refsdal 超新星的光由于引力透镜效应走上不同路径，从而在最初发现后近一年被再次探测到，相似的方法也可用于引力波的观测。\(^\text{[57]}\)尽管这项技术仍处于早期阶段，一种类似于手机通过 GPS 卫星三角定位（triangulation）确定位置的方法将帮助天文学家追踪引力波的起源。\(^\text{[58]}\)
\subsection{参见}
\begin{itemize}
  \item 引力波背景（Gravitational~wave~background）
  \item 引力波观测台（Gravitational\mbox{-}wave~observatory）
  \item 引力波观测列表（List~of~gravitational~wave~observations）
  \item 匹配滤波（Matched~filter）在引力波天文学中的应用（\#Gravitational\mbox{-}wave~astronomy）
\end{itemize}
\subsection{参考文献}
\begin{enumerate}
  \item Patrick~R.~Brady; Jolien~D.~E.~Creighton (2003)，《引力波天文学》（Gravitational~Wave~Astronomy），载于《物理科学与技术百科全书》（\textit{Encyclopedia~of~Physical~Science~and~Technology}），第~3~版，Academic~Press，页~33--48。
  \item Rabinarayan~Swain; Priyasmita~Panda; Hena~Priti~Lima; Bijayalaxmi~Kuanar; Biswajit~Dalai (2022~年~1--2~月)，《引力波：未来天文学综述》（Gravitational~waves:~A~review~on~the~future~astronomy），\textit{International~Journal~of~Multidisciplinary~Research~and~Growth~Evaluation}，第~3~卷第~1~期，页~38--50。
  \item Rothman,~Tony (2018~年~3~月)，《引力波的秘密历史》（The~Secret~History~of~Gravitational~Waves），American~Scientist。存档于~2024~年~3~月~20~日，检索于~2024~年~3~月~20~日。
  \item Moore,~Christopher; Cole,~Robert; Berry,~Christopher (2013~年~7~月~19~日)，《引力波探测器与源》（Gravitational~Wave~Detectors~and~Sources），检索于~2014~年~4~月~17~日。
  \item Overbye,~Dennis (2016~年~2~月~11~日)，《物理学家探测到引力波，证实爱因斯坦的理论》（Physicists~Detect~Gravitational~Waves,~Proving~Einstein~Right），New~York~Times。检索于~2016~年~2~月~11~日。
  \item Krauss,~Lawrence (2016~年~2~月~11~日)，《黑暗中的美》（Finding~Beauty~in~the~Darkness），New~York~Times。
  \item Pretorius,~Frans (2005)，《双黑洞时空的演化》（Evolution~of~Binary~Black-Hole~Spacetimes），Physical~Review~Letters}，95~(12):~121101。
  arXiv:gr-qc/0507014,~Bibcode:2005PhRvL..95l1101P,~doi:10.1103/PhysRevLett.95.121101。
  \item Campanelli,~M.; Lousto,~C.~O.; Marronetti,~P.; Zlochower,~Y. (2006)，《无切除的轨道黑洞双星精确演化》（Accurate~Evolutions~of~Orbiting~Black-Hole~Binaries~without~Excision），
  Physical~Review~Letters，96~(11):~111101。
 arXiv:gr-qc/0511048,~Bibcode:2006PhRvL..96k1101C,~doi:10.1103/PhysRevLett.96.111101。
  \item Baker,~John~G.; Centrella,~Joan; Choi,~Dae-Il; Koppitz,~Michael; van~Meter,~James (2006)，《从黑洞并合体系螺旋吸积配置中提取引力波》（Gravitational-Wave~Extraction~from~an~Inspiraling~Configuration~of~Merging~Black~Holes），
  Physical~Review~Letters，96~(11):~111102。
arXiv:gr-qc/0511103,~Bibcode:2006PhRvL..96k1102B,~doi:10.1103/PhysRevLett.96.111102。
  \item Abbott,~B.~P.~et~al. (2016~年~2~月~11~日)，《来自双黑洞并合的引力波观测》（Observation~of~Gravitational~Waves~from~a~Binary~Black~Hole~Merger），
  Physical~Review~Letters，116~(6):~061102。%
   arXiv:1602.03837,~Bibcode:2016PhRvL.116f1102A,~doi:10.1103/PhysRevLett.116.061102。
  \item Sesana,~A. (2013~年~5~月~22~日)，《在脉冲星定时频段中系统研究超大质量黑洞双星的预期引力波信号》（Systematic~Investigation~of~the~Expected~Gravitational~Wave~Signal~from~Supermassive~Black~Hole~Binaries~in~the~Pulsar~Timing~Band），Monthly~Notices~of~the~Royal~Astronomical~Society:~Letters，433~(1):~L1--L5。arXiv:1211.5375,~Bibcode:2013MNRAS.433L...1S,~doi:10.1093/mnrasl/slt034。
  \item IOPscience，《聚焦 NANOGrav 15 年数据集与引力波背景》（Focus~on~NANOGrav's~15~yr~Data~Set~and~the~Gravitational~Wave~Background）。
  \item “经过 15 年，脉冲星定时提供了宇宙引力波背景的证据”（After~15~years,~pulsar~timing~yields~evidence~of~cosmic~gravitational~wave~background），2023~年~6~月~29~日。
  \item O’Callaghan,~Jonathan (2023~年~6~月~28~日)，《巨大的引力“嗡鸣”在宇宙中传播》（An~Enormous~Gravity~‘Hum’~Moves~Through~the~Universe），Quanta~Magazine。
  \item Agazie,~Gabriella~et~al. (2023~年~6~月)，《NANOGrav 15 年数据集：引力波背景的证据》（The~NANOGrav~15~yr~Data~Set:~Evidence~for~a~Gravitational-wave~Background），The~Astrophysical~Journal~Letters}，951~(1):~L8。arXiv:2306.16213},~doi:10.3847/2041-8213/acdac6。
  \item Antoniadis,~J. (2023~年~6~月~28~日)，《欧洲脉冲星定时阵列第二次数据发布》（The~Second~Data~Release~from~the~European~Pulsar~Timing~Array），
  Astronomy~\&~Astrophysics}，678:~A50。
  arXiv:2306.16214},~doi:10.1051/0004-6361/202346844。
  \item Reardon,~Daniel~J.~et~al. (2023~年~6~月~29~日)，《利用帕克斯脉冲星定时阵列搜索各向同性引力波背景》（Search~for~an~Isotropic~Gravitational-wave~Background~with~the~Parkes~Pulsar~Timing~Array），
  The~Astrophysical~Journal~Letters}，951~(1):~L6。
  arXiv:2306.16215},~doi:10.3847/2041-8213/acdd02。
  \item Xu,~Heng~et~al. (2023~年~6~月~29~日)，《利用中国脉冲星定时阵列第一版数据搜索纳赫兹随机引力波背景》（Searching~for~the~Nano-Hertz~Stochastic~Gravitational~Wave~Background~with~the~Chinese~Pulsar~Timing~Array~Data~Release~I），
    Research~in~Astronomy~and~Astrophysics，23~(7):~075024。
    arXiv:2306.16216,~doi:10.1088/1674-4527/acdfa5。
  \item O’Callaghan,~Jonathan (2023~年~8~月~4~日)，《一种宇宙“嗡鸣”弥漫在宇宙中：科学家正竞相寻找其源头》（A~Background~‘Hum’~Pervades~the~Universe.~Scientists~Are~Racing~to~Find~Its~Source），
    Scientific~American，检索于~2023~年~8~月~5~日。
  \item “欧洲航天局研究‘不可见宇宙’的新愿景”（ESA’s~new~vision~to~study~the~invisible~universe），ESA，检索于~2013~年~11~月~29~日。
  \item Longair,~Malcolm (2012)，《宇宙世纪：天体物理与宇宙学史》（Cosmic~Century:~A~History~of~Astrophysics~and~Cosmology），剑桥大学出版社，ISBN~978-1-107-66936-9。
  \item Bahcall,~John~N. (1989)，《中微子天体物理学》（Neutrino~Astrophysics），剑桥大学出版社，ISBN~978-0-521-37975-5。
  \item Bahcall,~John (2000~年~6~月~9~日)，《太阳如何发光》（How~the~Sun~Shines），诺贝尔奖官方网站，检索于~2014~年~5~月~10~日。
  \item “1993~年诺贝尔物理学奖”（The~Nobel~Prize~in~Physics~1993），诺贝尔基金会，检索于~2014~年~5~月~3~日。
  \item Stairs,~Ingrid~H. (2003)，《利用脉冲星定时检验广义相对论》（Testing~General~Relativity~with~Pulsar~Timing），
    Living~Reviews~in~Relativity}，6~(1):~5。
    arXiv:astro-ph/0307536},~doi:10.12942/lrr-2003-5。
  \item Rincon,~Paul; Amos,~Jonathan (2017~年~10~月~3~日)，《爱因斯坦的“波”赢得诺贝尔奖》（Einstein’s~Waves~Win~Nobel~Prize），BBC~News。
  \item Overbye,~Dennis (2017~年~10~月~3~日)，《2017~年诺贝尔物理学奖授予 LIGO 黑洞研究团队》（2017~Nobel~Prize~in~Physics~Awarded~to~LIGO~Black~Hole~Researchers），New~York~Times。
  \item Kaiser,~David (2017~年~10~月~3~日)，《从引力波中学习》（Learning~from~Gravitational~Waves），New~York~Times。
  \item Nelemans,~Gijs~(2009~年~5~月~7~日)，《银河系引力波前景》（The~Galactic~Gravitational~Wave~Foreground），
        Classical~and~Quantum~Gravity，26~(9):~094030。
        arXiv:0901.1778,~Bibcode:2009CQGra..26i4030N,~doi:10.1088/0264-9381/26/9/094030。
  \item Stroeer,~A.; Vecchio,~A.~(2006~年~10~月~7~日)，《LISA 验证双星》（The~LISA~Verification~Binaries），
       lassical~and~Quantum~Gravity，23~(19):~S809–S817。
        arXiv:astro-ph/0605227,~Bibcode:2006CQGra..23S.809S,~doi:10.1088/0264-9381/23/19/S19。
  \item Abadie,~J.~et~al.~(2010~年~9~月~7~日)，《地基引力波探测器可观测的致密双星并合率预测》（Predictions~for~the~Rates~of~Compact~Binary~Coalescences~Observable~by~Ground-based~Gravitational-wave~Detectors），
        Classical~and~Quantum~Gravity，27~(17:~173001。
       rXiv:1003.2480~Bibcode:2010CQGra..27q3001A,~doi:10.1088/0264-9381/27/17/173001。
  \item “利用先进引力波探测器测量中等质量黑洞双星”（Measuring~Intermediate-Mass~Black-Hole~Binaries~with~Advanced~Gravitational~Wave~Detectors），
        伯明翰大学引力物理小组（University~of~Birmingham,~Gravitational~Physics~Group），检索于~2015~年~11~月~28~日。
  \item “观测中等质量黑洞的隐形碰撞”（Observing~the~Invisible~Collisions~of~Intermediate~Mass~Black~Holes），LIGO~科学合作组织（LIGO~Scientific~Collaboration），检索于~2015~年~11~月~28~日。
  \item Volonteri,~Marta; Haardt,~Francesco; Madau,~Piero~(2003~年~1~月~10~日)，《层级星系形成模型中超大质量黑洞的汇聚与并合历史》（The~Assembly~and~Merging~History~of~Supermassive~Black~Holes~in~Hierarchical~Models~of~Galaxy~Formation），
The~Astrophysical~Journal}，582~(2):~559–573。
arXiv:astro-ph/0207276},~Bibcode:2003ApJ...582..559V,~doi:10.1086/344675。
  \item Sesana,~A.; Vecchio,~A.; Colacino,~C.~N.~(2008~年~10~月~11~日)，《超大质量黑洞双星的随机引力波背景：对脉冲星定时阵列观测的启示》（The~Stochastic~Gravitational-wave~Background~from~Massive~Black~Hole~Binary~Systems:~Implications~for~Observations~with~Pulsar~Timing~Arrays），Monthly~Notices~of~the~Royal~Astronomical~Society，390~(1):~192–209。arXiv:0804.4476,~Bibcode:2008MNRAS.390..192S,~doi:10.1111/j.1365-2966.2008.13682.x。
  \item Amaro-Seoane,~Pau~et~al.~(2012~年~6~月~21~日)，《使用 eLISA/NGO 的低频引力波科学》（Low-frequency~Gravitational-wave~Science~with~eLISA/NGO），
   Classical~and~Quantum~Gravity，29~(12):~124016。
    arXiv:1202.0839,~Bibcode:2012CQGra..29l4016A,~doi:10.1088/0264-9381/29/12/124016。
  \item Amaro-Seoane,~P.~(2012~年~5~月)，《恒星动力学与极端质量比螺旋吸积》（Stellar~Dynamics~and~Extreme-mass~Ratio~Inspirals），Living~Reviews~in~Relativity}，21~(1):~4。arXiv:1205.5240},~Bibcode:2018LRR....21....4A,~doi:10.1007/s41114-018-0013-8。
  \item Berry,~C.~P.~L.; Gair,~J.~R.~(2012~年~12~月~12~日)，《利用引力波爆发观测银河系中心的超大质量黑洞》（Observing~the~Galaxy’s~Massive~Black~Hole~with~Gravitational~Wave~Bursts），Monthly~Notices~of~the~Royal~Astronomical~Society，429~(1):~589–612。arXiv:1210.2778,~Bibcode:2013MNRAS.429..589B,~doi:10.1093/mnras/sts360。
  \item Amaro-Seoane,~Pau~et~al.~(2007~年~9~月~7~日)，《中等与极端质量比螺旋吸积：天体物理、科学应用及 LISA 探测》（Intermediate~and~Extreme~Mass-ratio~Inspirals—Astrophysics,~Science~Applications~and~Detection~Using~LISA），
    Classical~and~Quantum~Gravity，24~(17):~R113–R169。
    arXiv:astro-ph/0703495,~Bibcode:2007CQGra..24R.113A,~doi:10.1088/0264-9381/24/17/R01。
  \item Gair,~Jonathan; Vallisneri,~Michele; Larson,~Shane~L.; Baker,~John~G.~(2013)，《使用低频空间引力波探测器检验广义相对论》（Testing~General~Relativity~with~Low-Frequency,~Space-Based~Gravitational-Wave~Detectors），
     Living~Reviews~in~Relativity，16~(1):~7。
    arXiv:1212.5575,~doi:10.12942/lrr-2013-7。
  \item Kotake,~Kei; Sato,~Katsuhiko; Takahashi,~Keitaro~(2006~年~4~月~1~日)，《核心坍缩超新星中的爆炸机制、中微子爆发与引力波》（Explosion~Mechanism,~Neutrino~Burst~and~Gravitational~Wave~in~Core-collapse~Supernovae），Reports~on~Progress~in~Physics，69~(4):~971–1143arXiv:astro-ph/0509456,~doi:10.1088/0034-4885/69/4/R03。
  \item Abbott,~B.~et~al.~(2007)，《来自未知孤立源与天蝎座~X-1~的周期性引力波搜索：LIGO 第二次科学运行结果》（Searches~for~Periodic~Gravitational~Waves~from~Unknown~Isolated~Sources~and~Scorpius~X-1:~Results~from~the~Second~LIGO~Science~Run），Physical~Review~D，76~(8):~082001。arXiv:gr-qc/0605028,~doi:10.1103/PhysRevD.76.082001。
  \item “在银河系中寻找最年轻的中子星”（Searching~for~the~Youngest~Neutron~Stars~in~the~Galaxy），LIGO~科学合作组织，检索于~2015~年~11~月~28~日。
  \item Binétruy,~Pierre~et~al.~(2012~年~6~月~13~日)，《宇宙引力波背景与 eLISA/NGO：相变、宇宙弦及其他源》（Cosmological~Backgrounds~of~Gravitational~Waves~and~eLISA/NGO:~Phase~Transitions,~Cosmic~Strings~and~Other~Sources），%
       Journal~of~Cosmology~and~Astroparticle~Physics，2012~(6):~027。
       arXiv:1201.0983,~doi:10.1088/1475-7516/2012/06/027。
  \item Damour,~Thibault; Vilenkin,~Alexander~(2005)，《来自宇宙（超）弦的引力辐射：脉冲、随机背景与观测窗口》（Gravitational~Radiation~from~Cosmic~(Super)strings:~Bursts,~Stochastic~Background,~and~Observational~Windows），
   Physical~Review~D}，71~(6):~063510。
     arXiv:hep-th/0410222},~doi:10.1103/PhysRevD.71.063510。
  \item Mack,~Katie~(2017~年~6~月~12~日)，《黑洞、宇宙碰撞与时空的涟漪》（Black~Holes,~Cosmic~Collisions~and~the~Rippling~of~Spacetime），%
     Scientific~American~(Blogs)。
  \item Schutz,~Bernard~F.~(2011~年~6~月~21~日)，《引力波探测器网络与三种评价指标》（Networks~of~Gravitational~Wave~Detectors~and~Three~Figures~of~Merit），
    Classical~and~Quantum~Gravity，28~(12):~125023。
    arXiv:1102.5421,~doi:10.1088/0264-9381/28/12/125023。
  \item Hu,~Wayne; White,~Martin~(1997)，《宇宙微波背景极化入门》（A~CMB~Polarization~Primer），
        New~Astronomy，2~(4):~323–344。
        arXiv:astro-ph/9706147,~doi:10.1016/S1384-1076(97)00022-5。
  \item Kamionkowski,~Marc; Stebbins,~Albert~(1997)，《宇宙微波背景极化统计》（Statistics~of~Cosmic~Microwave~Background~Polarization），
        Physical~Review~D，55~(12):~7368–7388。
        arXiv:astro-ph/9611125,~doi:10.1103/PhysRevD.55.7368。
  \item “规划光明的明天：利用先进 LIGO 与 Virgo 进行引力波天文学展望”（Planning~for~a~Bright~Tomorrow:~Prospects~for~Gravitational-wave~Astronomy~with~Advanced~LIGO~and~Advanced~Virgo），LIGO~科学合作组织，检索于~2015~年~12~月~31~日。
  \item Price,~Larry~(2015~年~9~月)，《寻找余辉：LIGO 视角》（Looking~for~the~Afterglow:~The~LIGO~Perspective），
    LIGO~Magazine~(7):~10。检索于~2015~年~11~月~28~日。
  \item 参见 Cutler~\&~Thorne~(2002)，第~2~节。
  \item 参见 Cutler~\&~Thorne~(2002)，第~3~节。
  \item 参见 Seifert~F.~et~al.~(2006)，第~5~节。
  \item 参见 Golm~\&~Potsdam~(2013)，第~4~节。
  \item “与爱因斯坦同行于曲径之上”（With~Einstein~on~Crooked~Paths）。
  \item “我们即将首次听到时空织物中的回声”（We~are~about~to~hear~echoes~in~the~fabric~of~space~for~the~first~time）。
  \item “引力透镜可能以史无前例的精度锁定黑洞并合事件”（Gravitational~Lenses~Could~Pin~Down~Black~Hole~Mergers~with~Unprecedented Accuracy）。
\end{enumerate}
\subsection{延伸阅读}
\begin{itemize}
\item Cutler, Curt; Thorne, Kip S. (2002), ``An overview of gravitational-wave sources'', 收录于 Bishop, Nigel; Maharaj, Sunil D.（编）,Proceedings of 16th International Conference on General Relativity and Gravitation (GR16), World Scientific, 第 4090 页, arXiv:gr-qc/0204090, Bibcode:2002gr.qc.....4090C, ISBN 978-981-238-171-2.中文译文：Cutler, Curt；Thorne, Kip S.（2002），“引力波源概述”，载于 Bishop, Nigel；Maharaj, Sunil D. （编），《第16届广义相对论与引力国际会议论文集（GR16）》，World Scientific 出版社，第 4090 页。
\item Thorne, Kip S. (1995), ``Gravitational radiation'', Particle and Nuclear Astrophysics and Cosmology in the Next Millennium: 160, arXiv:gr-qc/9506086,Bibcode:1995pnac.conf..160T.中文译文： Thorne, Kip S. （1995），“引力辐射”，《下一个千年的粒子与核天体物理学及宇宙学》，第 160 页。
\item Gravitational Wave Astronomy, Max Planck Institute for Gravitational Physics, 于 2013 年 2 月 6 日从原始版本存档，检索于 2013 年 1 月 24 日。中文译文：《引力波天文学》，马普引力物理研究所，2013 年 2 月 6 日存档，2013 年 1 月 24 日检索。
\item Schutz, B. F. (1999), ``Gravitational wave astronomy'', Classical and Quantum Gravity, 16(12A): A131 – A156, arXiv:gr-qc/9911034, Bibcode:1999CQGra..16A.131S, doi:10.1088/0264-9381/16/12A/307, S2CID 19021009.中文译文： Schutz, B. F. （1999），“引力波天文学”，《经典与量子引力》，第 16 卷 (12A)：A131 – A156 页。
\item LIGO Magazine, LIGO Scientific Collaboration.中文译文： 《LIGO 杂志》，LIGO 科学合作组织。
\item Seifert F.; Kwee P.; Heurs M.; Willke B.; Danzmann K. (2006), ``GEO600 Sensitivity'', 于 2009 年 12 月 23 日从原始版本存档。中文译文： Seifert F.；Kwee P.；Heurs M.；Willke B.；Danzmann K. （2006），“GEO600 灵敏度”，2009 年 12 月 23 日存档。
\end{itemize}
\subsection{外部链接}
\begin{itemize}
\item 
\textbf{中文译文：} LIGO 科学合作组织

\item \textbf{AstroGravS: Astrophysical Gravitational-Wave Sources Archive}\
\textbf{中文译文：} AstroGravS：天体物理引力波源档案

\item \textbf{Video (04:36) – Detecting a gravitational wave, Dennis Overbye, \textit{New York Times} (11 February 2016).}\
\textbf{中文译文：} 视频（时长 04:36）——《探测引力波》，作者 Dennis Overbye，发表于《纽约时报》（2016 年 2 月 11 日）。

\item \textbf{Video (71:29) – Press Conference announcing discovery: ``LIGO detects gravitational waves'', \textit{National Science Foundation} (11 February 2016).}\
\textbf{中文译文：} 视频（时长 71:29）——新闻发布会宣布发现：“LIGO 探测到引力波”，美国国家科学基金会（2016 年 2 月 11 日）。

\item \textbf{Gravitational Wave Astronomy}\
\textbf{中文译文：} 引力波天文学
\end{itemize}

---

是否希望我在下一步帮你将此部分与前面的“Further Reading”部分合并为一个完整的 \LaTeX 参考附录页（含统一排版、页眉页脚与参考编号）？
